\chapter{Further Reading}
\index{further reading}

This appendix collects books, standards documents, and online resources for continuing your study of Scheme and Kaappi.

\section{Books}

\begin{description}[leftmargin=1em, labelindent=0em]

\item[Structure and Interpretation of Computer Programs] \textit{Harold Abelson, Gerald Jay Sussman, and Julie Sussman} (MIT Press, 1996). The classic introduction to computer science using Scheme. Covers abstraction, recursion, interpreters, and compilation. Available free online from MIT Press.

\item[The Scheme Programming Language, 4th Edition] \textit{R.\ Kent Dybvig} (MIT Press, 2009). A thorough reference and tutorial for the Scheme language. Covers R6RS but most content applies to R7RS. Available free online at \texttt{scheme.com/tspl4/}.

\item[The Little Schemer] \textit{Daniel P.\ Friedman and Matthias Felleisen} (MIT Press, 1996). A short, playful introduction to recursive thinking using Scheme. Structured as a series of questions and answers that build up to writing a Scheme interpreter in Scheme. Highly recommended if you want to deepen your intuition for recursion.

\item[The Seasoned Schemer] \textit{Daniel P.\ Friedman and Matthias Felleisen} (MIT Press, 1996). The sequel to \textit{The Little Schemer}. Covers continuations, mutation, and advanced control flow.

\item[Essentials of Programming Languages, 3rd Edition] \textit{Daniel P.\ Friedman and Mitchell Wand} (MIT Press, 2008). Uses Scheme to build interpreters for progressively more complex languages. Excellent for understanding how programming languages work.

\end{description}

\section{Standards and Specifications}

\begin{description}[leftmargin=1em, labelindent=0em]

\item[R7RS-small Report] The official language standard that Kaappi implements. Defines all syntax forms, standard libraries, and semantics. Available at \texttt{r7rs.org}. Essential reading once you are comfortable with the language.

\item[SRFI Process] Scheme Requests for Implementation---the community process for proposing and standardizing library extensions. Browse all SRFIs at \texttt{srfi.schemers.org}. Each SRFI includes a rationale, specification, and reference implementation.

\end{description}

\section{Kaappi Resources}

\begin{description}[leftmargin=1em, labelindent=0em]

\item[Documentation Site] \texttt{kaappi-lang.org/docs/} --- the complete Kaappi reference, including the guided tutorial, language quick reference, and procedure documentation organized by category (600+ procedures).

\item[Tutorial] \texttt{kaappi-lang.org/docs/guide/tutorial/} --- a 12-section beginner tutorial that complements this book with interactive examples.

\item[Cookbook] \texttt{kaappi-lang.org/docs/cookbook/} --- practical recipes for common tasks: building REST APIs, parsing JSON, working with databases, writing CLI tools, and testing.

\item[Playground] \texttt{kaappi-lang.org/playground/} --- a browser-based REPL running Kaappi via WebAssembly. Useful for quick experiments without a local install.

\item[Source Code] \texttt{github.com/kaappi/kaappi} --- the Kaappi implementation in Zig. Reading the source is the best way to understand how the language works at a deep level.

\end{description}

\section{Online Resources}

\begin{description}[leftmargin=1em, labelindent=0em]

\item[Scheme Community Wiki] \texttt{community.schemewiki.org} --- tutorials, implementation comparisons, and community-maintained documentation.

\item[Awesome Scheme] A curated list of Scheme resources, libraries, and tools maintained on GitHub.

\item[Functional Programming Subreddit] \texttt{reddit.com/r/scheme} and \texttt{reddit.com/r/lisp} --- active communities for discussing Scheme and Lisp topics.

\end{description}
